\newpage
\part{Standard Configurations}

\section{NPX Guide}

\textbf{Node Package Execute (npx)} is commonly used to run pre-built MCP servers directly from the npm registry without global installation. 
This provides a "universal" configuration that works identically across AI agents and IDEs by using globally accessible commands.

\subsection{Essential Core Servers}
These servers represent the essential core for an AI agent's toolbelt, dividing capabilities into distinct pillars:

\begin{itemize}
    \item \textbf{Repository Control} (\lstinline[style=inline]|@modelcontextprotocol/server-github|): Allows any agent to fork, read files, check out branches, submit PRs, and review issues across codebases.
    \item \textbf{Smart Web Discovery} (\lstinline[style=inline]|@modelcontextprotocol/server-brave-search| or \lstinline[style=inline]|exa-mcp-server|): Essential for checking modern documentation, syntax changes, or software versions without hitting context limits.
    \item \textbf{Context Ingestion} (\lstinline[style=inline]|@modelcontextprotocol/server-fetch|): An official utility that visits any URL, strips out HTML clutter, and returns clean Markdown for the agent to process instantly.
    \item \textbf{Live System Auditing} (\lstinline[style=inline]|@modelcontextprotocol/server-postgres|): Gives the agent schema inspection and query-running execution capabilities directly against active application layers.
\end{itemize}

\subsection{Advanced Reasoning and Memory}
Adding these engines changes how your agent solves complex problems, moving away from blind session-by-session actions.

\begin{itemize}
    \item \textbf{Memory} (\lstinline[style=inline]|@modelcontextprotocol/server-memory|): Provides a local, graph-based knowledge storage system. It solves the amnesia problem between different context windows by letting the agent save and recall entities and structural rules persistently.
    \item \textbf{Sequential Thinking} (\lstinline[style=inline]|@modelcontextprotocol/server-sequential-thinking|): An official reasoning orchestrator designed to stop agents from rushing into incorrect code fixes. It forces the agent to explicitly state hypotheses, challenge assumptions, and track steps before touching workspace files.
\end{itemize}

\paragraph{Universal NPX Configuration Block.}
The following JSON block contains high-utility official distributions pulled straight from canonical registries. It can be placed in configuration files like \lstinline[style=inline]|claude_desktop_config.json|, \lstinline[style=inline]|windsurf_mcp_config.json|, or Gemini CLI's \lstinline[style=inline]|settings.json|.

\begin{lstlisting}[style=json]
{
  "mcpServers": {
    "github": {
      "command": "npx",
      "args": ["-y", "@modelcontextprotocol/server-github"],
      "env": {
        "GITHUB_PERSONAL_ACCESS_TOKEN": "YOUR_TOKEN_HERE"
      }
    },
    "brave-search": {
      "command": "npx",
      "args": ["-y", "@modelcontextprotocol/server-brave-search"],
      "env": {
        "BRAVE_API_KEY": "YOUR_API_KEY_HERE"
      }
    },
    "fetch-markdown": {
      "command": "npx",
      "args": ["-y", "@modelcontextprotocol/server-fetch"]
    },
    "memory": {
      "command": "npx",
      "args": ["-y", "@modelcontextprotocol/server-memory"]
    },
    "sequential-thinking": {
      "command": "npx",
      "args": ["-y", "@modelcontextprotocol/server-sequential-thinking"]
    }
  }
}
\end{lstlisting}
